\subsection{Giải pháp lưu trữ dữ liệu}

Hệ thống áp dụng mô hình Polyglot Persistence, sử dụng nhiều loại cơ sở dữ liệu khác nhau để tối ưu hóa cho từng loại dữ liệu và trường hợp sử dụng cụ thể.

\subsubsection{MongoDB - Document Database}

MongoDB là cơ sở dữ liệu NoSQL hướng tài liệu, được sử dụng để lưu trữ metadata và dữ liệu cấu hình của hệ thống Florus.

\textbf{Đặc điểm kỹ thuật:}
\begin{itemize}
    \item \textbf{Schema-less:} Cấu trúc dữ liệu linh hoạt, phù hợp với metadata đa dạng.
    \item \textbf{JSON-like Documents:} Lưu trữ dữ liệu dưới dạng BSON (Binary JSON).
    \item \textbf{Horizontal Scaling:} Hỗ trợ Sharding để mở rộng theo chiều ngang.
\end{itemize}

\textbf{Ứng dụng trong hệ thống:}
\begin{itemize}
    \item Lưu trữ thông tin User, Project, Permissions (Auth Service).
    \item Lưu trữ cấu hình các nguồn dữ liệu (Ingestion Service).
    \item Cache dữ liệu bảng để tăng tốc độ preview (Delta Service).
    \item Lưu trữ metadata của các mô hình (Model Service).
\end{itemize}

\subsubsection{Neo4j - Graph Database}

Neo4j là cơ sở dữ liệu đồ thị hàng đầu, được sử dụng để mô hình hóa và truy vấn các mối quan hệ phức tạp giữa người dùng và sản phẩm.

\textbf{Đặc điểm kỹ thuật:}
\begin{itemize}
    \item \textbf{Native Graph Storage:} Tối ưu hóa cho việc duyệt đồ thị (Graph Traversal).
    \item \textbf{Cypher Query Language:} Ngôn ngữ truy vấn trực quan cho đồ thị.
    \item \textbf{ACID Compliance:} Đảm bảo tính toàn vẹn dữ liệu.
\end{itemize}

\textbf{Ứng dụng trong hệ thống:}
\begin{itemize}
    \item Lưu trữ Knowledge Graph với các Node (User, Item, Category).
    \item Lưu trữ các Relationship (CLICKED, PURCHASED, BELONGS\_TO).
    \item Tính toán Graph Embeddings cho mô hình Recommendation.
    \item Hỗ trợ truy vấn Item-Item similarity và User-Item interactions.
\end{itemize}

\subsubsection{Redis - In-Memory Data Store}

Redis là cơ sở dữ liệu key-value trong bộ nhớ, được sử dụng làm bộ nhớ đệm tốc độ cao và lưu trữ Session State.

\textbf{Đặc điểm kỹ thuật:}
\begin{itemize}
    \item \textbf{Sub-millisecond Latency:} Thời gian truy xuất cực nhanh ($<1ms$).
    \item \textbf{Rich Data Structures:} Hỗ trợ String, List, Set, Sorted Set, Hash.
    \item \textbf{Redis Cluster:} Hỗ trợ sharding và replication cho High Availability.
\end{itemize}

\textbf{Ứng dụng trong hệ thống:}
\begin{itemize}
    \item \textbf{Session State:} Lưu chuỗi click của người dùng trong phiên hiện tại.
    \item \textbf{Graph Embeddings Cache:} Lưu vector đặc trưng được sync từ Neo4j.
    \item \textbf{Global Trending List:} Lưu Top-K sản phẩm phổ biến (Sorted Set).
    \item \textbf{Rate Limiting:} Kiểm soát số lượng request từ client.
\end{itemize}

\subsubsection{MinIO - Object Storage}

MinIO là hệ thống lưu trữ đối tượng hiệu năng cao, tương thích hoàn toàn với Amazon S3 API.

\textbf{Đặc điểm kỹ thuật:}
\begin{itemize}
    \item \textbf{S3-Compatible API:} Cho phép sử dụng các SDK/Tools của AWS S3.
    \item \textbf{Kubernetes-native:} Thiết kế tối ưu cho môi trường Container.
    \item \textbf{Erasure Coding:} Đảm bảo độ bền dữ liệu với chi phí lưu trữ thấp.
\end{itemize}

\textbf{Ứng dụng trong hệ thống:}
\begin{itemize}
    \item Lưu trữ MLflow Artifacts (Model weights, Metrics, Plots).
    \item Lưu trữ các file phi cấu trúc (Images, Videos, PDFs).
    \item Thay thế HDFS cho môi trường Cloud-native (được phân tích chi tiết tại Chương 3).
\end{itemize}

\subsubsection{So sánh và lựa chọn Database}

\begin{table}[H]
\centering
\caption{So sánh các giải pháp lưu trữ trong hệ thống}
\label{tab:database-comparison}
\begin{tabular}{|l|l|l|l|}
\hline
\textbf{Database} & \textbf{Loại dữ liệu} & \textbf{Latency} & \textbf{Use Case chính} \\
\hline
MongoDB & Document/JSON & 1-10ms & Metadata, Config \\
\hline
Neo4j & Graph & 10-100ms & Knowledge Graph \\
\hline
Redis & Key-Value & $<1ms$ & Cache, Session \\
\hline
MinIO & Object/Blob & 10-100ms & Files, Artifacts \\
\hline
\end{tabular}
\end{table}

Việc lựa chọn database phù hợp cho từng loại dữ liệu giúp tối ưu hóa hiệu năng tổng thể của hệ thống, đảm bảo đáp ứng yêu cầu về độ trễ thấp cho Serving Layer đồng thời duy trì tính nhất quán cho Processing Layer.
